\chapter{Spații euclidiene}

\section{Ortonormare și complement ortogonal}

Ne îndreptăm spre aplicațiile în geometria \emph{clasică}, adică
euclidiană, în care este esențial să putem calcula \emph{lungimi (distanțe)}
și \emph{unghiuri}.

În liceu, puteam face aceasta folosind produsul scalar al doi vectori
descompuși în reperul cartezian. Astfel, fie vectorii:
\[
  \vec{v} = a \vec{i} + b \vec{j}, \quad \vec{w} = c \vec{i} + d \vec{j}, %
  a, b, c, d \in \RR.
\]
Avem:
\begin{itemize}
\item \emph{produsul scalar:} $ \vec{v} \cdot \vec{w} = ac + bd \in \RR $;
\item \emph{lungimea vectorului} $ v = \sqrt{\vec{v} \cdot \vec{v}} $;
\item \emph{cosinusul unghiului între doi vectori:} %
  $ \cos(\widehat{\vec{v}, \vec{w}}) = \dfrac{\vec{v} \cdot \vec{w}}{v w} $;
\item vectorii sînt \emph{perpendiculari} dacă $ \vec{v} \cdot \vec{w} = 0 $.
\end{itemize}

Acestea sînt noțiunile pe care le extindem acum în cazul spațiilor vectoriale
arbitrare. Vom prezenta pentru cazul spațiului real $ V = \RR^n $, celelalte
rezolvîndu-se similar, prin intermediul izomorfismelor canonice.

Fie, deci, $ v = (v_1, \dots, v_n), w = (w_1, \dots, w_n) $ doi vectori
din $ \RR^n $. Se definesc:
\begin{itemize}\index{produs!scalar} \index{vector!norma} \index{vector!cosinus}
\item \textbf{produsul scalar:}
  \[
    \langle v, w \rangle = \sum_{i = 1}^n v_i w_i \in \RR;
  \]
\item \textbf{norma vectorului:}
  \[
    || v || = \sqrt{\langle v, v \rangle} \in \RR;
  \]
\item \textbf{cosinusul unghiului între doi vectori:}
  \[
      \cos(v, w) = \frac{\langle v, w \rangle}{||v|| \cdot ||w||} \in [-1, 1].
  \]
\end{itemize}

Avem și cazul particular de interes:
\begin{definition}\label{def:perp}
  Doi vectori $ v $ și $ w $ din $ \RR^n $ se numesc \emph{perpendiculari}
  (\emph{ortogonali}), notat $ v \perp w $, dacă $ \langle v, w \rangle = 0 $.
\end{definition}

Cu acestea, avem:
\begin{definition}\label{def:sp-eucl}
  \index{spații!euclidiene}
  Spațiul vectorial $ \RR^n $ se numește \emph{euclidian}, deoarece este
  înzestrat cu produsul scalar de mai sus.
\end{definition}

Aplicațiile care ne vor interesa sînt: calculul complementului
ortogonal al unui subspațiu și ortonormarea unei baze.

Începem cu primul subiect.

\begin{definition}\label{def:compl-ort}
  \index{spații!subspații!complement ortogonal}
  Fie $ V $ un subspațiu al unui spațiu euclidian ($ \RR^n $, în majoritatea
  aplicațiilor). Se definește:
  \[
    V^\perp = \{ x \in \RR^n \mid x \perp v, \forall v \in V \}
  \]
  subspațiu care se numește \emph{complementul ortogonal} al lui $ V $.
\end{definition}

O proprietate esențială, care ne va permite să găsim complementul ortogonal, este
următoarea:
\begin{theorem}\label{thm:compl-ort}
  Fie $ V $ un subspațiu al lui $ \RR^n $. Dacă $ V^\perp $ este complementul
  său ortogonal, atunci:
  \[
    V \oplus V^\perp \simeq \RR^n.
  \]
\end{theorem}

Rezultă, folosind teorema lui Grassmann și definiția sumei directe, că:
\begin{itemize}
\item $ \dim V^\perp = \dim \RR^n - \dim V $;
\item $ V^\perp \cap V = \{ 0 \} $.
\end{itemize}

De asemenea, amintim că în $ \RR^2 $, de exemplu, baza canonică formată
din $ \{(1, 0), (0, 1)\} $, pe care o putem înțelege și $ \{ \vec{i}, \vec{j} \} $,
are o proprietate specială. Mai precis, vectorii $ \vec{i} $ și $ \vec{j} $
sînt \emph{versori}, adică $ \vec{i} \perp \vec{j} $ și $ i = j = 1 $.

În orice spațiu euclidian, putem obține o bază similară, numită \emph{bază %
  ortonormată}, în două etape, folosind \textbf{procedeul Gram-Schmidt}.
\index{algoritm!Gram-Schmidt}

Fie, așadar, $ B = \{ b_1, \dots, b_n \} $ o bază arbitrară în $ \RR^n $.
Vrem să obținem din ea o bază \emph{ortonormată} $ C = \{ c_1, \dots, c_n \} $,
adică să fie formată din versori:
\begin{itemize}
\item $ c_i \perp c_j, \forall i \neq j $;
\item $ || c_i || = 1, \forall i $.
\end{itemize}

Procedeul construiește pas cu pas $ c_i $, pornind de la $ b_i $.
\begin{enumerate}[(1)]
\item Alegem $ c_1 = b_1 $;
\item Definim $ c_2 = b_2 + \alpha c_1 $, cu $ \alpha $ un scalar pe care vrem
  să-l determinăm. Punem condiția $ c_2 \perp c_1 $ și rezultă:
  \[
    0 = \langle c_2, c_1 \rangle = \langle b_2 + \alpha c_1, c_1 \rangle = %
    \langle b_2, c_1 \rangle + \alpha \langle c_1, c_1 \rangle.
  \]
  Așadar:
  \[
    \alpha = - \frac{\langle b_2, c_1 \rangle}{\langle c_1, c_1 \rangle}.
  \]
\item Continuăm mai departe să definim $ c_3 = b_3 + \alpha c_2 + \beta c_1 $
  și determinăm $ \alpha, \beta $ din condițiile $ c_3 \perp c_2 $ și $ c_3 \perp c_1 $.
\item Procedeul continuă după regula generală:
  \[
    c_i = b_i + \sum_{j = i - 1}^1 \alpha_j c_j,
  \]
  iar condițiile pe care le vom putea folosi la fiecare pas, pentru determinarea
  constantelor $ \alpha_j $ sînt $ c_i \perp c_j, \forall j < i $, deoarece $ c_j $
  au fost determinate la pașii anteriori.
\end{enumerate}

Rezultă, cu aceasta, o bază \emph{ortogonală} $ C = \{ c_1, \dots, c_n \} $.
Această bază se \emph{normează}, obținînd vectori de normă $ 1 $ foarte simplu.
Definim:
\[
  c'_i = \frac{1}{||c_i||} \cdot c_i
\]
și observăm că acum avem $ ||c'_i|| = 1 $.

Așadar, baza $ C' = \{ c'_1, \dots, c'_n \} $ este o \emph{bază ortonormată},
adică formată doar din versori, iar procedeul se încheie.

%%%%%%%%%%%%%%%%%%%%%%%%%%%%%%%%%%%%%%%%%%%%%%%%%%%%%%%%%%%%%%%%%%%%%%
\section{Aplicații în geometrie}

Amintim, înainte de generalizarea într-un spațiu vectorial arbitrar, că
în geometria analitică din liceu am studiat și condiții de coliniaritate
pentru doi vectori. Această noțiune poate fi generalizată ușor:

\begin{definition}\label{def:vectori-col}
  \index{spații!euclidiene!vectori coliniari}
  Fie $ v, w \in \RR^n $ doi vectori, care se scriu într-o bază fixată
  $ B = \{ b_i \} $ a lui $ \RR^n $ cu coordonatele $ \{ \alpha_i \} $,
  respectiv $ \{ \beta_i \} $.

  Spunem că cei doi vectori sînt coliniari, notat $ v \parallel w $
  dacă și numai dacă au coordonatele proporționale:
  \[
    \dfrac{\alpha_1}{\beta_1} = \dfrac{\alpha_2}{\beta_2} = \dots = %
    \dfrac{\alpha_n}{\beta_n} = k \in \RR.
  \]
  Evident, dacă $ k = 1 $, avem $ v = w $.
\end{definition}

De asemenea, în liceu, dacă este dat un vector din plan
$ \vec{v} = a \vec{i} + b \vec{j} $, lui i se poate calcula \emph{proiecția}
pe axa OX calculînd produsul scala $ \vec{v} \cdot \vec{i} = a $.

Noțiunea poate fi generalizată astfel:

\begin{definition}\label{def:proiectie-subspatiu}
  Fie $ v \in \RR^n $ un vector arbitrar și $ V $ un subspațiu al lui $ \RR^n $
  de bază $ B = \{ b_1, \dots, b_t \} $.

  \emph{Proiecția ortogonală a vectorului $ v $ pe subspațiul $ V $} se definește
  prin:
  \[
    v_{\parallel V} = \dr{pr}_V v = \sum_{i = 1}^t \dfrac{\langle v, b_i\rangle}{%
      \langle b_i, b_i \rangle} b_i.
  \]
\end{definition}

Geometric, $ \dr{pr}_V v $ este \qq{cel mai apropiat} de $ v $ vector din $ V $.

Mai mult, un rezultat important este:
\begin{theorem}\label{thm:proi-perp}
  În condițiile și cu notațiile de mai sus, $ v = v_{\parallel V} + v_{\perp V} $,
  unde $ v_{\perp V} \in V^\perp $.
\end{theorem}

Cîteva aplicații geometrice simple urmează.

În plan, produsul scalar poate fi folosit pentru a calcula aria paralelogramului.

\begin{proposition}\label{pr:aria-paral}
  Produsul scalar al vectorilor $ u, v \in \RR^2 $ este egal cu aria
  paralelogramului determinat de cei doi vectori.
\end{proposition}

În spațiu, fie doi vectori:
\begin{align*}
  \vec{a} &= a_1 \vec{i} + a_2 \vec{j} + a_3 \vec{k} \\
  \vec{b} &= b_1 \vec{i} + b_2 \vec{j} + b_3 \vec{k}.
\end{align*}

\index{produs!vectorial}
\emph{Produsul vectorial} al celor doi vectori, notat $ \vec{a} \times \vec{b} $
poate fi calculat ca un determinant formal:
\[
  \vec{a} \times \vec{b} = %
  \begin{vmatrix}
    \vec{i} & \vec{j} & \vec{k} \\
    a_1 & a_2 & a_3 \\
    b_1 & b_2 & b_3
  \end{vmatrix}.
\]

Geometric, produsul vectorial produce un vector care este perpendicular pe
planul determinat de cei doi factori, a cărui orientare este dată cu regula
burghiului.

\begin{proposition}\label{pr:vol-paral}
  Norma (modulul) produsului vectorial este egală cu volumul paralelipipedului
  determinat de cei trei vectori, ca muchii adiacente.
\end{proposition}

De asemenea, avem și:
\begin{proposition}\label{pr:prod-vec-paral}
  Păstrînd notațiile și condițiile de mai sus,
  \[
    \vec{a} \parallel \vec{b} \Leftrightarrow \vec{a} \times \vec{b} = \vec{0}.
  \]
\end{proposition}

\index{produs!mixt}
Totodată, se poate defini și \emph{produsul mixt} pentru trei vectori.

Fie cei doi vectori de mai sus și $ \vec{c} = c_1 \vec{i} + c_2 \vec{j} + c_3 \vec{k} $>

Produsul mixt al celor trei vectori se calculează prin:
\[
  \vec{a} \cdot (\vec{b} \times \vec{c}) = %
  \begin{vmatrix}
    a_1 & a_2 & a_3 \\
    b_1 & b_2 & b_3 \\
    c_1 & c_2 & c_3
  \end{vmatrix}.
\]

\begin{proposition}\label{pr:vectori-copl}
  Păstrînd notațiile și contextul de mai sus, vectorii $ \vec{a}, \vec{b} $
  și $ \vec{c} $ sînt coplanari dacă și numai dacă produsul lor mixt este nul.
\end{proposition}

\begin{proposition}\label{pr:volum-tetra}
  Păstrînd notațiile și contextul de mai sus, volumul tetraedrului determinat
  de vectorii $ \vec{a}, \vec{b}, \vec{c} $ ca muchii adiacente se calculează:
  \[
    V = \dfrac{1}{6} \vec{a} \cdot (\vec{b} \times \vec{c}).
  \]
\end{proposition}


%%%%%%%%%%%%%%%%%%%%%%%%%%%%%%%%%%%%%%%%%%%%%%%%%%%%%%%%%%%%%%%%%%%%%%

\section{Exerciții}

1. Să se ortonormeze bazele:
\begin{enumerate}[(a)]
\item $ B_1 = \{ (1, 0, 1), (-1, 2, 0), (0, 1, -1) \} $ a lui $ \RR^3 $;
\item $ B_2 = \{ X + 1, X^2 + 1, 3 \} $ a lui $ \RR_2[X] $;
\item $ B_3 = \{ (1, -1, 0, 0), (0, 1, 1, -1), (1, 0, -1, 1), (2, 1, 0, 1) \} $
  a lui $ \RR^4 $.
\end{enumerate}

\vspace{0.5cm}

2. Să se completeze mulțimea $ \{ (1, -1, 2), (1, 2, 3) \} $ pînă la o bază
a lui $ \RR^3 $ și să se ortonormeze baza rezultată.

\vspace{0.5cm}

3. Să se completeze mulțimea $ \{ 3, 2X + 1 \} $ pînă la o bază a lui $ \RR_2[X] $
și să se ortonormeze baza rezultată.

\vspace{0.5cm}

4. Găsiți complementele ortogonale pentru:
\begin{enumerate}[(a)]
\item $ V_1 = \dr{Sp} \{ (-1, 1, 2), (0, 1, 1) \} $ în $ \RR^3 $;
\item $ V_2 = \dr{Sp} \{ (1, 1, 1) \} $ în $ \RR^3 $;
\item $ V_3 = \{ (x, y, z) \in \RR^3 \mid 3x + y = 0 \} $ în $ \RR^3 $;
\item $ V_4 = \{ (x, y, z) \in \RR^3 \mid x = y \} $ în $ \RR^3 $;
\item $ V_5 = \{ p \in \RR_2[X] \mid p'(0) = 0 \} $ în $\RR_2[X] $;
\item $ V_6 = \{ p \in \RR_2[X] \mid p(0) = 0 \} $ în $ \RR_2[X] $;
\item $ V_7 = \dr{Sp}\{ X \} $ în $ \RR_2[X] $;
\item $ V_8 = \dr{Sp} \{ 2, 5X^2 \} $ în $ \RR_2[X] $;
\item $ V_9 = \{ (x, y, z, t) \in \RR^4 \mid 3x + y = 0, y + z - t = 0 \} $
  în $ \RR^4 $.
\end{enumerate}

\vspace{0.5cm}

5. Folosind produsele scalare canonice, calculați $ \langle u, v \rangle, ||u||, ||v|| $
și proiecțiile $ \dr{pr}_u v, \dr{pr}_v u $. Calculați și cosinusul unghiului între
$ u $ și $ v $ și decideți dacă vectorii sînt ortogonali:
\begin{enumerate}[(a)]
\item $ u = (1, 2), v = (-1, 2) \in \RR^2 $;
\item $ u = (1, 1, 1), v = (1, -2, 0) \in \RR^3 $;
\item $ u = 1 + X, v = X^2 \in \RR_2[X] $;
\item $ u = \begin{pmatrix} 1 & 0 \\ 2 & 1 \end{pmatrix} $,
  $ v = \begin{pmatrix} 0 & -1 \\ 1 & 0 \end{pmatrix} \in M_2(\RR) $;
\end{enumerate}

\vspace{0.5cm}

6. Fie subspațiul $ W = \dr{Sp} \{ v_1 = (1, 0, 1, 1), v_2 = (1, -1, 1, 0) \} \seq \RR^4 $.
\begin{enumerate}[(a)]
\item Determinați $ W^\perp $;
\item Arătați că $ W \oplus W^\perp \simeq \RR^4 $;
\item Pentru $ v = (1, 1, 1, 1) $, determinați $ v_0 = \dr{pr}_{v_1} v $;
\item Pentru $ v $ de mai sus, determinați $ v_{\parallel W} v \in W $
  și $ v_{\perp W} = v - v_{\parallel W} $.
\end{enumerate}

\vspace{0.5cm}

7. Aflați proiecția $ v_{\parallel W} $ a vectorului $ v $ pe
subspațiul $ W $, precum și componenta sa ortogonală $ v_{\perp W} \in W^\perp $ pentru:
\begin{enumerate}[(a)]
\item (*) $ v = 1 + x \in \RR_2[x], W = \dr{Sp}\{ p_1 = 1 + x^2, p_2 = 1 \} $, cu
  produsul scalar definit prin:
  \[
    \langle p, q \rangle = \int_{-1}^1 p(t) q(t) dt, \quad \forall p, q \in \RR_2[x].
  \]
\item $ v = (1, 2, 1), W = \dr{Sp} \{ v_1 = (2, 1, 0), v_2 = (-1, 4, 1) \} \seq \RR^3 $;
\item $ v = \begin{pmatrix} 1 & 2 \\ 4 & 1 \end{pmatrix} $, $ W = \dr{Sp} %
    \left\{ C = \begin{pmatrix} 1 & 0 \\ 0 & 1 \end{pmatrix}, %
      D = \begin{pmatrix} 0 & 1 \\ 2 & 1 \end{pmatrix} \right\} \seq M_2(\RR) $;
\item $ v = (2, 1, -1), W = \{ (x, y, z) \mid x + y - 2z = 0 \} \seq \RR^3 $.
\end{enumerate}
%%% Local Variables:
%%% mode: latex
%%% TeX-master: "../ag-20-21"
%%% End:
